% Transcription — fonds Alexandre Grothendieck, shelfmark 161-6, batch 2 (pages 21-33).
% Established from the facsimile archives/batches/161-6/batch-02.pdf (a mirror of
% https://grothendieck.umontpellier.fr/161-6.pdf), per the skill
% transcribe-grothendieck. The batch file carries no cover sheet: PDF page N is
% archivists' page 20 + N, checked against the pencilled numbers.
%
% Pass: Opus 5.5 (claude-opus-5-5), 2026-09-22 — first pass, unchecked against the pages by a human.
%
% What the batch holds, in the archivists' order. Loose leaves, several hands
% of paper (yellow ruled, white ruled), no pagination of his own; the batch
% is three separate subjects, not one argument.
%
%   21     Icosahedra over a base S: a quadratic form Q, "icosaèdres locaux"
%          with cos'(D_i, D_j) = c, c^2 + c - 1 = 0; statements a)-c)
%          (existence, local conjugacy under SO(Q), isotropy group finite
%          étale with geometric fibres A_5); the computation cos α = 1/√5,
%          cos 2α = -3/5, whence B(x,y)^2 = (4/5) Q(x)Q(y), and with
%          B' = B/2, B'(x,y)^2 = (1/5) Q(x)Q(y). A small diagram of
%          μ_5, μ_5*, V_5, Ic_Q, S, Z[1/4 (sic?)] at top left.
%   22     Written the other way up and cancelled by a large cross, but
%          legible: for a topos E and a small category C, π(C,E) =
%          Hom(C°,E), a 2-functor Cat × Top → Top, and the comparison
%          morphism γ(C,E): π(C,E) → Ĉ × E; question whether it is always an
%          equivalence, with cases a) sums, b) C a groupoid, c) E = D̂.
%   23     Blank: absent.
%   24-30  A run on isotopy classes of tame embeddings φ: X → S of a compact
%          1-dimensional triangulable X into an oriented surface S: the
%          cyclic orders ω_x(φ) at the vertices, Ω(X) = ∏ Ω_x(X) with
%          (ν_x - 1)! elements each; Th 1 (the tubular neighbourhood U_α
%          depends only on α); U_α - X = V_α a union of half-open annuli;
%          the data (α, π, g, ν) and the model surfaces S_{α,π,g,ν}; Th 2
%          (every embedding is isomorphic to a model one); reduction to
%          Isot(S_ρ, S) (a torsor under Autext π_1) and to the stabiliser
%          H of the restriction to X; the Teichmüller group; the "normal"
%          case π_1(X) → π_1(S) surjective, with plane (S^2) embeddings
%          classified by the α with γ(α) = 0; a question on H (p. 29) and a
%          fibration sequence G/H ≃ X marked "??" (p. 30, rest blank).
%   31     Blank: absent.
%   32-33  Graphs as (S, A, A' → S × S, σ); G = group of symmetries,
%          transitive on vertices and edges; case a) all vertices isolated,
%          Aut Γ = Perm(S) × Aut(A, σ); case b) (cancelled by two strokes,
%          legible) all edges "circular", A ≃ G/H, S ≃ G/H', σ̇ in
%          (H' ∩ Norm_G(H))/H of order 2; case c) (p. 33) edges not circular,
%          G ⊃ H' ⊃ H, σ̇ ∈ Norm_G(H)/H, faithfulness ∩ g^{-1}Hg = {1}, and
%          the condition H' ∩ σ̇^{-1}H'σ̇ ⊂ H for Γ to be simplicial.
%
% Notation fixed for the batch. Page 21: his double-struck letters are
% \boldsymbol{\mu} and \mathbb{V}; the golden-ratio letter is written γ
% (it is a γ or a looped r). Page 22: underlined Hom is
% \underline{\mathrm{Hom}}, C° is C^{\circ}, the presheaf topos \hat{C}.
% Pages 24-30: U_α and V_α in his ronde are \mathcal{U}_\alpha,
% \mathcal{V}_\alpha, Pl_isot is \mathrm{Pl}_{\mathrm{isot}}; the number of
% ends is written with a letter read β on p. 26 and ν on p. 27 (see the
% notes), both kept as written; on p. 30 the three script letters are
% \mathcal{G}, \mathcal{H}, \mathfrak{X}. Pages 32-33: σ̇ is \dot{\sigma}.
%
% Doubtful, not settled: Z[1/4] and the index "5s" on p. 21; the
% connective prose of pp. 25 and 28 (interlinear insertions, the "*"
% margin note of p. 28); the exponent of S^2 on p. 28, written like a 1 and
% read 2 from the argument; the labels under φ and L on p. 32. On p. 33,
% H' and H are glossed G_a and G_s, the reverse of p. 32; left as written.
%
% For the sibling batch and the modernised reading. Pages 24-30 are the
% combinatorial description of ribbon graphs (cyclic orders at vertices,
% the thickening U_α, boundary cycles, gluing surfaces of genus g_λ with
% ν_λ punctures) — the "graphes" of the folder title, well before the
% Esquisse. Page 22 has nothing to do with the rest of the batch. The
% dating line is the inventory's own, for the shelfmark and its group.
\documentclass[11pt,a4paper]{article}
% Le préambule commun à toutes les transcriptions du fonds Grothendieck.
%
% Il tient en une idée : **l'incertitude se compose, elle ne se résout pas.**
% Une transcription qui lisse un mot douteux a détruit la seule chose pour
% laquelle on la faisait — un lecteur doit pouvoir savoir, à chaque endroit,
% ce qui était sur le feuillet et ce qui a été ajouté. Les six macros qui
% suivent sont donc l'appareil critique, et non de la décoration.
%
% Les mêmes commandes sont reconnues par `scripts/render.mjs`, qui produit la
% vue de lecture du site. Une macro ajoutée ici doit l'être aussi là-bas,
% faute de quoi le rendu échouera bruyamment — ce qui est voulu : un rendu
% silencieusement faux est pire qu'un rendu absent.

\ProvidesPackage{grothendieck}[2026/08/08 Transcriptions du fonds Grothendieck]

\usepackage{amsmath,amssymb,amsthm}
\usepackage{mathrsfs}
\usepackage{stmaryrd}
% Les diagrammes commutatifs sont partout dans ce fonds : c'est la langue
% même de la géométrie algébrique de Grothendieck, et une transcription qui
% les rendrait en prose ne transcrirait rien.
\usepackage{tikz-cd}
% Français par défaut — c'est la langue des feuillets — mais l'anglais est
% chargé aussi : la traduction et le résumé sont des documents anglais, et la
% typographie française (espaces avant les deux-points) y serait une faute.
% Ces documents basculent par \selectlanguage{english} en tête de corps.
\usepackage[english,french]{babel}
\usepackage{fontspec}
\usepackage{geometry}
\geometry{a4paper,margin=25mm,marginparwidth=28mm}
\usepackage{marginnote}
\usepackage{xcolor}
% ulem, not soul. `soul`'s \st and \ul are built on a character-by-character
% scan that XeLaTeX's font machinery breaks: the document fails with an
% undefined control sequence rather than degrading. ulem does the same two jobs
% and is Unicode-engine safe. `normalem` stops it hijacking \emph, which the
% transcriptions use for Grothendieck's own emphasis.
\usepackage[normalem]{ulem}
\usepackage{hyperref}
\hypersetup{colorlinks=true,linkcolor=black,urlcolor=black,pdfborder={0 0 0}}

% --- Métadonnées du lot -----------------------------------------------------
% Reprises telles quelles par le script de rendu, qui les affiche en tête de
% la vue de lecture. Elles fixent ce que le fichier prétend couvrir : sans
% elles, une transcription flotte sans point d'attache dans un fonds de seize
% mille feuillets.

\newcommand{\folder}[1]{\gdef\grothFolder{#1}}
\newcommand{\batch}[1]{\gdef\grothBatch{#1}}
\newcommand{\leaves}[2]{\gdef\grothFirst{#1}\gdef\grothLast{#2}}
\newcommand{\foldertitle}[1]{\gdef\grothFoldertitle{#1}}
\gdef\grothFolder{?}\gdef\grothBatch{?}\gdef\grothFirst{?}\gdef\grothLast{?}\gdef\grothFoldertitle{}

% --- Le résumé --------------------------------------------------------------
%
% Il ouvre la lecture modernisée, sous le titre « Résumé », et il est composé
% en petit. Ce n'est pas un ornement : c'est le seul passage du document qui
% ne soit pas des mathématiques, et un lecteur doit pouvoir le reconnaître
% avant de l'avoir lu — puis le sauter s'il connaît déjà le sujet, ou s'y
% arrêter s'il ne le connaît pas. Le filet qui le suit marque la reprise du
% corps.

\newenvironment{resume}
  {\small\setlength{\parskip}{0.45em}}
  {\par\medskip\hrule\medskip}

% --- Mots-clés ---------------------------------------------------------------
%
% En anglais, en clôture du résumé de la lecture modernisée : le vocabulaire
% moderne sous lequel chercher ce que les feuillets construisent. Le manifeste
% du site (`npm run manifest`) les extrait pour étiqueter la cote entière —
% cette ligne est la source unique des tags, il n'y a pas de fichier à côté.
\newcommand{\keywords}[1]{\par\smallskip\noindent
  {\footnotesize\textsc{Keywords} --- \textit{#1}}\par}

% --- L'appareil critique ----------------------------------------------------

% Le numéro de feuillet, en marge. C'est l'ancre qui permet de régler un
% désaccord sur une formule en regardant *un* feuillet plutôt qu'une chemise
% de six cents. Le site s'en sert aussi pour tourner les pages du fac-similé
% quand on fait défiler la transcription.
\newcommand{\leaf}[1]{%
  \marginnote{\footnotesize\color{gray!70}\textsf{#1}}%
  \label{leaf:#1}\ignorespaces}

% Illisible : on ne devine pas. Un mot inventé qui se lit comme les autres est
% le pire résultat possible.
\newcommand{\ill}{\textcolor{red!60!black}{[\,\ldots\,]}}

% Lecture douteuse : le mot est donné, et signalé comme incertain.
\newcommand{\uncertain}[1]{\uline{#1}}

% Ajout de l'éditeur — une lettre manquante, un mot sous-entendu.
\newcommand{\add}[1]{\textcolor{blue!50!black}{[#1]}}

% Biffé par Grothendieck lui-même. Ce qu'il a rayé fait partie du document :
% une rature dit souvent où la pensée a changé de direction.
\newcommand{\struck}[1]{\sout{#1}}

% Note du transcripteur, sur l'état matériel du feuillet.
\newcommand{\note}[1]{\par\smallskip{\small\color{gray!60!black}\textit{[#1]}}\par\smallskip}

% Note marginale de Grothendieck, distincte du corps du texte.
\newcommand{\marginal}[1]{\par\smallskip\noindent\hspace{1em}%
  {\small\color{gray!60!black}#1}\par\smallskip}

% --- Titre ------------------------------------------------------------------

\AtBeginDocument{%
  \begin{center}
    {\large\bfseries Fonds Alexandre Grothendieck --- cote n\textsuperscript{o}~\grothFolder}\\[2pt]
    {\itshape\grothFoldertitle}\\[4pt]
    {\small feuillets \grothFirst--\grothLast\ (lot \grothBatch)}\\[2pt]
    {\footnotesize\color{gray!60!black}Transcription établie d'après le fac-similé numérisé
     par l'Université de Montpellier.\\
     Les lectures incertaines sont soulignées, les illisibles notées [\,\ldots\,],
     les ajouts entre crochets.}
  \end{center}
  \vspace{1em}}

\endinput


\folder{161-6}
\batch{2}
\pages{21}{33}
\dating{[à partir de 1973-vers 1977] — le groupe « Dossiers rassemblés par Grothendieck sur différentes thématiques » (161-1 à 162-6) est daté [après 1961-vers 1977]}
\watermark{Édition de démonstration}
\foldertitle{Graphes, icosaèdre [etc.] : notes manuscrites (s.d.)}

\begin{document}

\page{21}

\note{en haut à gauche, une colonne de symboles reliés par des traits
verticaux sans flèche, et à côté un petit diagramme fléché ; on les donne
tels quels, sans décider ce que les traits signifient}

\[
\begin{array}{c}
\boldsymbol{\mu}_5 \\ \vert \\ \boldsymbol{\mu}_5^{*} \\ \vert \\ \mathbb{V}_5 \\ \vert \\ \mathbb{Z}[\tfrac{1}{4}]
\end{array}
\]
\note{la dernière ligne se lit $\mathbb{Z}[\tfrac14]$ ; on attendrait
$\mathbb{Z}[\tfrac15]$, mais le chiffre a la forme d'un 4. Les lettres
$\mathbb{V}$ sont tracées doublées, comme $\boldsymbol{\mu}$}

\begin{tikzcd}
 & \mathrm{Ic}_{Q} \arrow[d] \\
 & \mathbb{V}_{5s} \arrow[d] \\
\mathbb{V}_5 & S \arrow[l] \arrow[u, bend left=40]
\end{tikzcd}

\note{l'indice $5s$ de $\mathbb{V}_{5s}$ est une lecture douteuse ; la
flèche courbe remonte de $S$ vers $\mathbb{V}_{5s}$, à côté de la flèche
descendante. $\mathbb{V}_5$ est celui de la colonne de gauche}

\[
\cos'(D_i, D_j)^2 + \cos'(D_i, D_j)
\]

$\cos'(D_i, D_j)$ indépendant de $i, j$ pour $i \neq j$, soit $c$

$c$

\begin{itemize}
  \item[a)] $\exists$ des icos. loc. \marginal{$\uncertain{q} = c$}
  \ill{} tels que $c$ donné \marginal{$c^2 + c - 1 = 0$}
  \item[b)] Deux icos. sont loc. conjugués par $SO(Q)$
  \item[c)] Le groupe d'isotropie d'un icosaèdre dans $SO(Q)$ (groupe
  des automorphismes d'un icosaèdre) est fini étale $/S$, et ses fibres
  géom. isom. au groupe \emph{alterné} à 5 \emph{lettres}.
\end{itemize}
\note{« dans $SO(Q)$ » et « fini » sont des ajouts interlinéaires de sa
main ; « $/S$ » est une lecture douteuse}

\note{dans la marge gauche : un icosaèdre dessiné en perspective ; plus bas
un cercle traversé de deux diamètres presque confondus, puis un cercle avec
un diamètre prolongé hors du cercle et deux points marqués sur le bord}

icosaèdre projectif --- icosaèdre

\[
1 - \gamma = \gamma^2, \qquad 2 - \gamma = 1 + \gamma^2, \qquad
\gamma^2 + \gamma - 1 = 0
\]
\note{la lettre notée ici $\gamma$ est un $\gamma$ ou un $r$ bouclé ; elle
est le nombre d'or inverse, racine de $\gamma^2+\gamma-1=0$, comme le $c$
plus haut}

\[
\cos\alpha = \frac{\struck{4}\gamma/2}{1 - \gamma/2} = \frac{\gamma}{2-\gamma}
= a + b\gamma = \frac{1}{5}(2\gamma + 1) = \frac{1}{\sqrt{5}}
\]
\note{la dernière égalité est entourée ; un signe biffé la précède}

\[
\begin{aligned}
\cos 2\alpha &= \cos^2\alpha - \sin^2\alpha = 2\cos^2\alpha - 1
  = 2\,\frac{\gamma^2}{4 - 4\gamma + \gamma^2} - 1 \\
&= \frac{2(1-\gamma)}{4 - 4\gamma + 1 - \gamma} - 1
  = \frac{2 - 2\gamma}{5(1-\gamma)} - 1 = \frac{2}{5} - 1 = -\frac{3}{5}
\end{aligned}
\]
\[
\frac{2}{25}\,\underbrace{[4\gamma^2 + 4\gamma + 1]}_{4 - 4\gamma + 4\gamma + 1} - 1
\]

\[
\frac{B(x,y)^2}{Q(x)Q(y)} - 2 = -\frac{6}{5}
\qquad\qquad 2\cos 2\alpha = -\frac{6}{5}
\]

\[
\boxed{B(x,y)^2 = \frac{4}{5}\,Q(x)Q(y)}
\]

\[
2a + (2b - 1 \uncertain{-a})\gamma + b\gamma^2 = 0, \qquad
\gamma^2 + \Bigl(-2 + \frac{1}{b} + \frac{a}{b}\Bigr)\gamma - \frac{2a}{b} = 0
\]
\[
-\frac{2a}{b} = -1, \quad \boxed{b = 2a}, \qquad
-2 + \frac{1}{2a} + \frac{1}{2} = 1, \qquad a = \frac{1}{5},\ b = \frac{2}{5}
\]

en car. 2, \uncertain{on fait} $B(x,y)^2 = 0$

\note{les deux lignes qui suivent sont d'une encre plus pâle}
\[
\cos' 2\alpha = 2\cos 2\alpha = 2(2\cos^2\alpha - 1) = (2\cos\alpha)^2 - 2,
\qquad \frac{1}{2a} = \frac{5}{2}, \qquad a = \frac{1}{5}
\]

donc si l'on \ill{} \uncertain{$\cos'^2$} \ill{}

ceci signifie : \struck{$y \perp x$} $y \parallel x$

\[
\tfrac{1}{2}B = B' \qquad\qquad
\boxed{B'(x,y)^2 = \frac{1}{5}\,Q(x)Q(y)}
\]
\note{sous « $\frac12 B = B'$ », trois segments issus d'un même point}

\page{22}

\note{la page, écrite dans l'autre sens de la feuille, est barrée de deux
longs traits en croix ; elle se lit entièrement dessous}

$E$ topos

$C$ petite catégorie

alors $\underline{\mathrm{Hom}}(C^{\circ}, E) = \pi(C, E)$ un topos

Si on a un morphisme $C \xrightarrow{f} C'$, $g : E \to E'$ d'où
$g^{*} : E' \to E$, on conclut un foncteur $\pi(f, g) : \pi(C, E) \to \pi(C', E')$
\[
\pi(f, g)^{*} = (u' \mapsto g^{*} \circ u' \circ f^{\circ}), \qquad
\pi(f, g)_{*} = (u' \mapsto g_{*} f_{*}^{(E)})
\]
\marginal{vérifier l'expression \uncertain{donnée} \uncertain{en} $\pi(f, g)_{*}$ …}
qui est le foncteur image inverse associé à un morphisme de topos
\[
\pi(f, g) : \pi(C, E) \longrightarrow \pi(\uncertain{C'}, E')
\]
\note{la première lettre de $\pi(C', E')$ ressemble ici à un $E$ ; une
ligne plus haut la même flèche est écrite $\pi(C,E) \to \pi(C',E')$}

de sorte que $\pi(C, E)$ dépend de façon covariante de $(C, E)$ [i.e. est
un 2-foncteur $(\mathrm{Cat}) \times (\mathrm{Top}) \to (\mathrm{Top})$]

On va définir un morphisme de topos \struck{fonctoriel} en $C, E$
\[
\text{(1)}\qquad
\boxed{\gamma = \gamma(C, E) : \pi(C, E) \longrightarrow \hat{C} \times_{\mathrm{top}} E}
\]
en définissant les composantes
\[
\begin{aligned}
\gamma_1 &= \gamma_1(C, E) : \pi(C, E) \to \hat{C} \\
\gamma_2 &= \gamma_2(C, E) : \pi(C, E) \to E
\end{aligned}
\]
par
\[
\gamma_1^{*} : \hat{C} \to \pi(C, E) \quad \text{défini par} \quad
F \mapsto [X \mapsto F(X)_E]
\]
où $\hat{C} = \pi(C, P)$, $P$ topos ponctuel ; i.e.
$\gamma_1 = \pi(\mathrm{id}_C, p_E)$, où $p_E : E \to P$ « unique »
morphisme de topos ;
\[
\gamma_2^{*} : E \to \pi(C, E), \qquad X \mapsto (i \mapsto X)
\]
où $E \simeq \pi(C_0, \uncertain{E})$, $C_0$ catégorie ponctuelle ;
$\gamma_2^{*}(X)$ est le foncteur constant de valeur $X$, i.e.
$\gamma_2 = \pi(p_C, \mathrm{id}_E)$, $p_C : C \to C_0$ unique morphisme
\uncertain{de topos}.

\textbf{Question} $\gamma(C, E)$ est-il \uncertain{toujours} une équivalence ?

\begin{itemize}
  \item[a)] \struck{Exemple} \struck{Supposons} On note que $\pi(C, E)$
  transforme sommes en un des deux arguments en sommes, et de même pour
  $\hat{C} \times E$ ($C \mapsto \hat{C}$ transformant sommes en sommes).
  Donc on est ramené à traiter le cas $C$ connexe.
  \item[b)] $C$ un groupoïde --- \uncertain{fini} par supposition connexe.
  Alors $\hat{C} \simeq B_G$, $G$ un groupe, $\pi(C, E) \simeq B_{G_E}$
  (catégorie des objets de $E$ avec $G$ opérant), ok
  \item[c)] $E = \hat{D}$, alors $\pi(C, E) \cong \widehat{C \times D}$, et
  $\gamma(C, E)$ est le morph. de topos \struck{fonctoriel} évident
  $\widehat{C \times D} \to \hat{C} \times \hat{D}$
\end{itemize}

\page{24}

$X$ espace topologique \struck{compact} $C_1$ [triangulable compact de
dim $\leq 1$]

$S$ surface \emph{orientée}

Nous voulons étudier les classes d'isotopie des plongements (modérés …)
$X \xrightarrow{\varphi} S$

Pour tout tel plongement $\varphi$, et pour tout sommet \struck{de}
\uncertain{non isolé} topologique de $X$, il y a un \emph{ordre}
\struck{\uncertain{circulaire}} $\omega_x(\varphi)$ sur l'ens. des arêtes
(orientation) issues de $x$ induit par $\varphi$ \struck{de} [donnée vide si
le sommet $x$ est un bout, i.e. s'il n'y a qu'une seule arête issue de $x$ ;
sinon, il y \struck{\ill{}} a $\nu_x \geq 3$ telles arêtes, et \struck{\ill{}}
l'\uncertain{ens.} des ordres circulaires possibles est un ens.
$\Omega_x(X)$ à $(\nu_x - 1)!$ éléments].
\note{« topologique » est ajouté au-dessus de la ligne ; « $\omega_x(\varphi)$ »
aussi, avec un trait qui le rattache à « ordre »}

On trouve ainsi un invariant
\[
\omega(\varphi) = (\omega_x(\varphi))_{x \in X'_0} \in \prod_{x \in X'_0} \Omega_x(X)
\]
\note{sous $X'_0$ : « ens. des sommets de $X$ qui ne sont ni isolés ni des
bouts »}

On a une appl.
\[
\mathrm{Pl}_{\mathrm{isot}}(X, S) \xrightarrow{\ \omega\ } \prod_{x \in X'_0} \Omega_x(X)
\overset{\mathrm{df}}{=} \Omega(X)
\]
et notre propos est d'étudier les fibres \struck{de} $\omega$,
$\mathrm{Pl}^{\alpha}_{\mathrm{isot}}(X, S)$, plus généralement, pour $X$ et
un $\alpha = (\alpha_x)_{x \in X'_0} \in \struck{\ill{}}\ \Omega(X)$ fixés,
d'étudier pour la $S$ les plongements $\varphi : X \hookrightarrow S$ tels que
$\omega(\varphi) = \alpha$.

Pour ceci, \struck{\ill{}} \struck{D}éfinissons pour le plongement « le »
voisinage tubulaire \struck{\uncertain{fidèle}} \uncertain{fermé} de celui-ci. On
trouve

\textbf{Th 1} Pour $\alpha$ fixé, ce voisinage tubulaire \struck{est
\ill{}} ne dépend, à homéomorphismes près induisant l'identité sur $X$ [de
façon unique à isotopie près] que de $\alpha$, soit $\mathcal{U}_{\alpha}$, et
non de $S$. \struck{Donc} On trouve donc une application
\[
\mathrm{Pl}^{\alpha}_{\mathrm{isot}}(X, S) \longrightarrow
\mathrm{Pl}_{\mathrm{isot}}(\mathcal{U}_{\alpha}, S)
\]
Celle-ci est une \uncertain{bijection}

\page{25}

Ce théorème ramène essentiellement l'étude de $\mathrm{Pl}^{\alpha}_{\mathrm{isot}}$
à celle de $\mathrm{Pl}_{\mathrm{isot}}(\mathcal{U}_{\alpha}, S)$. Pour faire
cette étude, on doit pouvoir donner un procédé pour décrire combinatoirement
$\mathcal{U}_{\alpha}$ en termes de $\alpha$. Notons

\textbf{\emph{Prop}} $\mathcal{U}_{\alpha} - X = \mathcal{V}_{\alpha}$ a des
composantes connexes $\mathcal{V}^{i}_{\alpha}$ qui sont
\struck{le bord de $\mathcal{V}_{\alpha}$ est} des produits
$S^1 \times \struck{\ill{}}\,[1, 2[$ \struck{et \ill{}} \struck{le filtre des}
avec $\mathrm{bord}(\mathcal{U}_{\alpha}) \cap \mathcal{V}^{i}_{\alpha} \simeq
S^1 \times \{1\}$, le filtre des voisinages de $X$ dans $\mathcal{U}_{\alpha}$
induisant sur $\mathcal{V}^{i}_{\alpha}$ le filtre des voisinages de
$S^1 \times \{1\}$ dans $S^1 \times\, ]0, 2[$.
\note{$S^1$ est glosé « cercle » ; la borne de $[1,2[$ et de $]0,2[$ est
récrite sur un autre chiffre, et le passage biffé de la deuxième ligne est
lourdement raturé. Les indices $\{1\}$ et $]0,2[$ sont donnés tels
qu'écrits}

\marginal{le nb des $\mathcal{V}^{i}_{\alpha}$ est $b_1(X) + 1$ si $X$
connexe non vide ? Cela donne en général $b_1(X) + b_0(X)$}

\struck{Donc}

Donc la surface $S$, pour un plongement donné de $\mathcal{U}_{\alpha}$, peut
se décrire comme \struck{$\mathcal{U}_{\alpha}$} déduite de $\mathcal{U}_{\alpha}$
en lui collant une variété orientée à bord $\mathcal{U}'$ [compacte ssi $S$
est compacte] --- savoir $\overline{S - \varphi(\mathcal{U}_{\alpha})}$ --- à
$\mathcal{U}_{\alpha}$, par un isom. de $\partial\,\mathcal{U}'$ avec
$-\partial\,\mathcal{U}_{\alpha} \simeq$
\note{la ligne s'arrête sur « $\simeq$ »}

\note{à gauche, plusieurs croquis : une grande courbe fermée contenant deux
disques accolés, une boucle et une petite figure ramifiée ; plus bas, un
anneau (deux courbes fermées emboîtées) avec une petite boucle attachée au
bord intérieur ; en marge, un autre anneau, où le disque intérieur porte
\uncertain{« dans le »}}

\emph{Le Pb} de trouver $\mathrm{Pl}^{\alpha}_{\mathrm{isot}}(S)$ \ill{} \ill{}
i.e. de trouver, à \struck{\uncertain{homéom.}} \uncertain{isom.} près
\add{induisant l'identité sur $X$}, tous les plongements de $X$ dans une
surface \struck{\ill{}} orientée, on peut déjà envisager la solution de ce
Pb : il s'agit de trouver les façons possibles d'écrire
\[
I \times S^1 \qquad (I = \pi_0(\mathcal{V}_{\alpha}))
\]
($I$ étant l'ens. d'indices fini discret) comme espace des bords orienté
d'une surface à bord [connexe ou non]
\struck{[NB $\overline{M}$ si $\mathcal{U}'_{\alpha}$ \ill{} connexe, \ill{}
composantes connexes de \ill{} $\mathcal{U}'_{\alpha}$ \ill{} rencontrent le
bord \ill{} alors $S$ \ill{} (\ill{})]} $S$ est connexe.
\note{« induisant l'identité sur $X$ » est une lecture de l'ajout
interlinéaire, qui est serré et en partie douteux ; « $= \pi_0(\mathcal{V}_\alpha)$ »
est écrit au-dessus de $I$}

Donc la classification est donnée aussi $\{$\uncertain{relativement}$\}$ aux
$\mathcal{U}'$ quant à leur \ill{}

\marginal{NB $\mathcal{U}'_{\alpha}$ non nécessairement connexe, mais on peut
supposer que ses composantes connexes ont un bord (\uncertain{en tt cas}
toute composante connexe de $S$ rencontre $X$)}

\page{26}

\begin{itemize}
  \item[a)] On donne une partition $\pi$ de $I$ en parties
  $(I_{\lambda})_{\lambda \in \Lambda}$
  \item[b)] Pour tout $\lambda$, on donne une \struck{variété} surface orientée
  connexe à bord $\mathcal{U}'_{\lambda}$, et un isom.
  $\partial\,\mathcal{U}'_{\lambda} \simeq -I_{\lambda} \times S^1$.
\end{itemize}
\note{« connexe » est ajouté au-dessus de la ligne}

Si on se borne aux $S$ « algébriques » (compactes moins ens. fini), une telle
$\mathcal{U}'_{\lambda}$ sera connue quand on connaît
\begin{itemize}
  \item[a)] Le genre $g_{\lambda}$ de la surface compacte sans bord obtenue en
  « rajoutant les points à l'infini » (i.e. les bouts) et en contractant
  chaque composante du bord en un point
  \item[b)] Le nb de bouts $\beta_{\lambda}$ de $\mathcal{U}'_{\lambda}$.
\end{itemize}

[\emph{NB} Comme il y aura tjrs un homéom. orienté de $\mathcal{U}'_{\lambda}$
sur elle-même \struck{\ill{}, \ill{} isotope à l'identité,} qui échange entre
eux de façon prescrite les éléments du bord, on n'a pas à s'en faire pour
préciser la correspondance entre \struck{\ill{}} composantes du bord et les
\struck{éléments de $I_{\lambda}$}]
\note{« orienté » et « composantes » sont ajoutés au-dessus de la ligne}

En fait, en termes de $g_{\lambda}$, $\beta_{\lambda}$, on construit
$\mathcal{U}'_{\lambda}$ en prenant la surface \struck{\ill{}} sans bord
orientée compacte de genre $g$, soit $S_g$, on y prend $I_{\lambda}$
rondelles disjointes $D_j$ ($j \in I_{\lambda}$), et \struck{\ill{}}
\struck{$\beta_{\lambda}$ points \ill{}} \uncertain{un ens.} fini $P_{\lambda}$
de cardinal $\beta_{\lambda}$ \uncertain{non rencontré} par
$\bigcup_{j \in I_{\lambda}} D_j$, et prend
\[
\mathcal{U}'_{\lambda} = S_g - \bigcup_{j \in I_{\lambda}} \mathring{D}_j - P_{\lambda}
\]

\marginal{\uncertain{Sauf si} \ill{} les $\mathcal{U}'_{\lambda}$
\uncertain{compactes}, i.e. $\beta_{\lambda} = 0$}

Soit $S_{\alpha, \pi, g, \beta}$ la surface ainsi construite à l'aide de
\[
\begin{cases}
\alpha \in \Omega(X) & [\text{d'où } I = I(\alpha), \text{ à expliciter combinatoirement}] \\
\pi \text{ partition de } I \text{ en } (I_{\lambda})_{\lambda \in \Lambda} & \\
(g_{\lambda})_{\lambda \in \Lambda} \in \mathbb{N}^{\Lambda} & \\
(\beta_{\lambda})_{\lambda \in \Lambda} \in \mathbb{N}^{\Lambda} &
\end{cases}
\]
\note{devant les deux $\mathbb{N}^{\Lambda}$, un premier symbole est
lourdement biffé}

On a donc un plongement canonique
\[
X \xrightarrow{\ \varphi_{\alpha, \pi, g, \beta}\ } S_{\alpha, \pi, g, \beta}
\]

\marginal{\emph{NB} $S_{\alpha, \pi, g, \beta}$ a $\beta = \sum_{\lambda}
\beta_{\lambda}$ bouts. Le genre calculé \uncertain{en} genre, il doit être
\uncertain{quelque chose} $\sum_{\lambda \in \Lambda} g_{\lambda} +
\gamma(\alpha)$, $\gamma(\alpha)$ un entier (le genre de
$\mathcal{U}_{\alpha}$) --- \emph{vérifier}}
\note{la note marginale est reliée par un trait à l'accolade des quatre
données ; la lettre du nombre de bouts y est la même que dans le texte, lue
$\beta$ partout ici, mais elle pourrait être un $\nu$ (cf. $\nu_x$ page~24)}

On trouve donc (S surfaces orientées \emph{algébriques})

\textbf{\emph{Th 2}} Tout plongement $X \to S$ est isomorphe à un
$\varphi_{\alpha, \pi, g, \beta}$ pour $\alpha, \pi, g, \beta$ convenables,
uniquement déterminés …
\note{la phrase s'arrête au bas de la page}

\page{27}

Donc on trouve, pour $S$ fixé, une application
\[
\mathrm{Pl}_{\mathrm{isot}}(X, S) \longrightarrow
\text{l'ensemble des systèmes } \alpha, \pi, g, \nu \text{ associés à } X
\]
(à expliciter combinatoirement)
\note{ici la lettre du nombre de bouts est nettement un $\nu$ ; c'est celle
que la page~26 écrit d'un trait qu'on a lu $\beta$}

et on est donc ramené à étudier, pour $\rho = \alpha, \pi, g, \nu$ fixé,
l'ens. des $\mathrm{Pl}_{\mathrm{isot}}(X, S; \rho)$ qui correspondent à
$\rho$ (non vide ssi $S \simeq S_{\rho} \overset{\mathrm{df}}{=}
S_{\alpha, \pi, g, \nu}$). Or on a une application \struck{Isot}
« restriction à $X$ »
\[
(*)\qquad
\mathrm{Isot}(S_{\rho}, S) \xrightarrow{\ r\ } \mathrm{Pl}_{\mathrm{isot}}(X, S; \rho)
\]
qui, d'après ce qui vient d'être vu, est \emph{surjective}.
\note{l'appel $(*)$ est le nôtre : la page encadre la flèche entre deux
barres verticales, et c'est à elle que renvoie plus bas « les deux membres
de $(*)$ ». Au-dessus de « ce qui vient », un ajout interlinéaire
\uncertain{« et précisé … »}}

\marginal{\uncertain{ce qui va suivre donnera} …}

\struck{Il reste :}

\begin{itemize}
  \item[a)] étudier quelles sont les fibres de cette application, i.e.
  \item[b)] Déterminer l'ens. $\mathrm{Isot}(S_{\rho}, S)$
\end{itemize}

La question b) est classique et résolue : $\mathrm{Isot}(S_{\rho}, S)$ est un
torseur à droite sous $\mathrm{Autext}(\pi_1(S_{\rho}))$ (ou à g. sous
$\mathrm{Autext}(\pi_1(S))$) [du moins si $S_{\rho}$ est compacte i.e.
$\sum_{\lambda \in \Lambda} \nu_{\lambda} = 0$ i.e. les $\nu_{\lambda}$ nuls
…]. \struck{Donc la question a) se ramène à la suivante :}
\note{« à droite » et « (ou à g. sous $\mathrm{Autext}(\pi_1(S))$) » sont
ajoutés au-dessus de la ligne. Deux longs traits obliques traversent la
suite de la page ; ils ne biffent pas le texte, qui est repris page~28}

Soit $f : S_{\rho} \xrightarrow{\sim} S_{\rho}$ un homéomorphisme, dont la
classe d'homotopie est déterminée par l'\struck{\ill{}} extérieur induit par
$f$ sur $\pi_1(S_{\rho}) = G_{\rho}$, soit $E(f) \in \mathrm{Autext}(G_{\rho})$,
\struck{soit $E(f)$} considérons $f \vert X = f_X$. Quelle relation y a-t-il
entre $E(f)$, $E(g)$ assurant que $f_X \overset{\mathrm{isot}}{\simeq} g_X$ ?

\marginal{\emph{NB}}

\begin{tikzcd}
G_{\rho} \arrow[r, bend left=20, "\bar{f}"] \arrow[r, bend right=20, "\bar{g}"'] & G_{\rho} \\
\pi_1(X) \arrow[u] \arrow[ur, no head, bend right=30] &
\end{tikzcd}

\note{croquis en marge gauche, rattaché à « NB » par un trait : deux flèches
parallèles $G_\rho \rightrightarrows G_\rho$ marquées $\bar f, \bar g$, une
flèche montante de $\pi_1(X)$ vers le premier $G_\rho$, et une courbe sans
tête de $\pi_1(X)$ au second}

Notons que le groupe
\[
G = \mathrm{Isot}_{o}(S, S) \qquad [\simeq G_{\rho} = \mathrm{Isot}_{o}(S_{\rho}, S_{\rho})]
\]
opère sur les deux membres de $(*)$, et \uncertain{c'est clair} $r$ commutant
à l'opération de ce groupe. D'autre part, \struck{comme} \uncertain{pour} le
premier membre est un torseur sous $G$, donc, choisissant
$f : S_{\alpha} \simeq S$, \struck{si} $r(f) = f_X = f \vert X$ \ill{} son
image, et $H$ son stabilisateur, on a que $r$ s'identifie à l'application
$G \to G/H$. Donc on

\page{28}

et finalement \uncertain{on arrive} aux 2 questions suivantes
\begin{itemize}
  \item[a)] Déterminer $G = \mathrm{Isot}(S, S)$
  \item[b)] Déterminer le sous-groupe $H$ de $G$ formé des $h : S \to S$ tels
  que $h \circ f \vert X \overset{\mathrm{isot}}{\simeq} f$.
\end{itemize}

On peut évidemment supposer $S = S_{\rho}$ donc $G = G_{\rho}$, et
$f = \varphi_{\rho} : X \to S_{\rho}$, dans la question dernière
\begin{itemize}
  \item[a)] Déterminer $G_{\rho} = \mathrm{Isot}(S_{\rho}, S_{\rho})$ \quad
  $S_{\rho}$ \uncertain{surface de genre} \ill{} moins
  \item[b)] Déterminer le sous-groupe $H_{\rho} \subset G_{\rho}$ des
  $h \in G_{\rho}$ tels que \struck{$h \vert X \simeq$}
  $h \circ \varphi_{\rho} \simeq \varphi_{\rho}$
\end{itemize}
\note{au-dessus de b), un ajout interlinéaire \uncertain{« $\nu'$, $\rho^k$ … »}
reste obscur}

Tout ceci \uncertain{semble} \ill{} un $X$ compact ($\in \mathrm{Ob}\, C_1$)
\uncertain{plutôt que} d'un couple $X \subset S_{\rho}$, $S_{\rho}$ surface
orientée avec \struck{\ill{}} (connexe)

La réponse à a) est classique (du moins si $S_{\rho}$ compacte) :
\struck{\ill{}} L'homomorphisme canon.
\[
G \longrightarrow \mathrm{Autext}(\pi_1(S, s))
\]
est un isom., et $G$ est donc le groupe de \emph{Teichmüller} de \struck{\ill{}}
$S$. Reste à trouver (pour un $X \underset{\varphi}{\subset} S$ plongé dans
$S$) quels sont les $h \in G$ tels que $h\varphi \underset{\mathrm{isot}}{\simeq}
\varphi$.

\emph{Cas particulier} ($X$ connexe et) $\pi_1(X) \to \pi_1(S)$ surjectif
(sauf erreur, ça signifie que la partition $\pi$ de $I$ est discrète
($\Lambda = I$) et que les $g_{\lambda}$ et les $\nu_{\lambda}$ sont nuls ---
i.e. que $S_{\rho}$ est déduit de $\rho$ simplement en « bouchant
\uncertain{simplement} \ill{} par des rondelles ») : plongements
\emph{normaux}*. Alors
\[
h\varphi \underset{\mathrm{isot}}{\simeq} \varphi \;\Longrightarrow\;
h\varphi \underset{\mathrm{homot}}{\simeq} \varphi
\]
donc

\begin{tikzcd}
\pi_1(X) \arrow[r, hook, "c"] \arrow[dr] & \pi_1(S) \arrow[d, "\pi_1(h)"] \\
 & \pi_1(S)
\end{tikzcd}

\note{sous $\pi_1(X)$ un trait vertical court, sans tête ni but, à côté de
la diagonale}

\[
\pi_1(h) \circ c \simeq c \quad \text{mod automorphismes intérieurs}
\]
donc $\pi_1(h)$ est un automorphisme intérieur, $c$ étant surjectif, en gros.
\uncertain{si} $X$ \emph{connexe}

\marginal{* \ill{} \uncertain{faisant} \ill{} \uncertain{préciser} \ill{}}

Cela nous donne en particulier les plongements dans \uncertain{$S^2$} (car tt
\struck{plongement} homéom. orienté de \uncertain{$S^2$} sur lui-même est isotope à
l'identité).
\note{l'exposant de $S^2$, aux deux occurrences, a la forme d'un 1 ; c'est
la sphère que l'argument demande, le cercle n'ayant pas de sens ici.
« homéom. orienté » est écrit au-dessus du mot biffé}

Ils sont classifiés par les $\alpha$ tels que $\gamma(\alpha)$ [le genre de
$\tilde{\mathcal{U}}_{\alpha}$ déduit de $\mathcal{U}_{\alpha}$ par le procédé
dit qui donne \struck{\ill{}} un

\page{29}

plongement normal$) = 0$. Il faut donner une certaine \uncertain{condition}
simple sur $\alpha$ pour qu'il en soit ainsi.
\note{la phrase commencée au bas de la page~28 se ferme ici :
« $\gamma(\alpha)$ [le genre de $\tilde{\mathcal{U}}_\alpha$ …] $= 0$ »}

\textbf{\emph{Question}} Supposons que $h \in G$ soit tel que pour toute
composante connexe $X_i$ de $X$, désignant par $\varphi_i : X_i \to S$ le
plongement induit,
\[
\pi_1(h \varphi_i) = \pi_1(h) \circ \pi_1(\varphi_i) \in
\struck{\ill{}}\ \mathrm{Homext}(\pi_1(X_i), \pi_1(S))
\]
soit égal à $\pi_1(\varphi_i)$. \struck{Alors a-t-on} (condition évidemment
nécessaire pour que $h \circ \varphi$ soit isotope à $\varphi$, \struck{donc}
i.e. $h \in H$) ; alors a-t-on $h \in H$ ?
\note{la question est tenue par un trait vertical en marge gauche. Dans
$\pi_1(h\varphi_i)$, un $\varphi$ surchargé précède le $h$ ; « Homext » est
écrit au-dessus d'un mot biffé}

\page{30}

\[
\mathcal{G}/\mathcal{H} \simeq \mathfrak{X} \qquad
\pi_1(\mathcal{H}) \to \pi_1(\mathcal{G}) \to \pi_1(\mathfrak{X}) \to
\pi_0(\mathcal{H}) \to \pi_0(\mathcal{G}) \to \pi_0(\mathfrak{X})
\]
\note{une longue accolade marquée « ?? » court au-dessus de la suite, de
$\pi_1(\mathcal{G})$ jusque vers $\pi_0(\mathfrak{X})$, où elle finit en
flèche. Les lettres $\mathcal{G}$, $\mathcal{H}$, $\mathfrak{X}$ rendent
une ronde, une cursive et un $X$ à double jambage}

\struck{$\pi_1(\mathcal{G}) \to \pi_1(\mathcal{H})$}
\uncertain{les homéom.} $h$ tels que $h \vert X = \text{identité}$
\[
\mathcal{H} = \mathrm{Aut}(S / X)
\]
\uncertain{le groupe top.}
\note{la ligne sur les $h$ tels que $h\vert X = \mathrm{id}$ est reliée
par un trait à $\mathcal{H}$ ; le reste de la page est blanc}

\page{32}

\struck{Graphes $(S, A, R, \sigma)$ [$S$ sommets, $A$ arêtes, $R$ relation
d'incidence orientée, $\sigma$ symétrie] $[R \subset S \times A]$}
\note{ce premier essai est encadré et hachuré}

\struck{Graphe $S$ : ensemble $S$ avec partie $A$ de $S \times S$, stable par
symétrie}

\[
\text{Graphes}\quad
(S, A, A' \xrightarrow{\ \varphi = (\varphi_1, \varphi_2) = (\varphi_1, \varphi_1 \circ \sigma)\ } S \times S, \sigma)
\;\;\langle\simeq\rangle\;\;
(S, R \subset A \times A, L \to A, \sigma)
\]
\note{sous la formule, des légendes fléchées : $S$ « sommets », $A$
« arêtes orientées », sous $\varphi$ un mot surchargé \ill{}, $\sigma$
« symétrie dans $A$ » ; une accolade renvoie à « [arêtes non circulaires
isolées] » ; $R$ « stable par sym. », $L$ « arêtes \uncertain{bouclées} »}

$G$ groupe des symétries du graphe \qquad
$G \simeq \mathrm{Perm}(S) \times \mathrm{Perm}(A)$

Supposons que $G$ est transitif sur $A$ et sur $S$.

Alors \struck{\ill{}} ou bien tous les \struck{points} sommets sont isolés (et
les arêtes non circulaires isolées) dans le graphe et partant formé de deux
ensembles $S$, $A$, $A' = \emptyset$, et un automorphisme d'ordre 2 dans $A$
\struck{\ill{}} sans point fixe ; \struck{l'action de $G$. Donc}
\[
\mathrm{Aut}\,\Gamma = \mathrm{Perm}(S) \times \underbrace{\mathrm{Aut}(A, \sigma)}
\]
\marginal{pour les sommets isolés}
Si $A_1 = A/\sigma$, on trouve que $\mathrm{Aut}(A, \sigma) \simeq
\mathrm{Perm}(A_1) \cdot (\mathbb{Z}/2\mathbb{Z})^{A_1}$ (produit
semi-direct)

\note{la suite de la page, jusqu'au crochet final, est traversée de deux
longs traits obliques parallèles ; elle se lit dessous}

[b) ou bien \struck{tous les sommets sont} il y a des sommets non isolés,
\struck{\ill{}}. Dans ce cas, il y a des arêtes non circulaires isolées
\struck{\ill{}} toutes les arêtes sont non isolées circulaires [car il y en
a non isolées]. \struck{Il est possible} Ou bien toutes les arêtes sont
circulaires, dans \uncertain{ce cas} $\varphi$ se réduit à une application
\emph{surjective}
\[
\varphi : A \longrightarrow S
\]
compatible avec $\sigma$ $[\varphi\sigma x = \varphi x]$
\marginal{il existe des arêtes non isolées \uncertain{circulaires}}

\struck{\ill{} $A$} Choisissons $a \in A$, donc
$A \simeq G/G_a^{=H}$ (groupe d'isotropie) et si $s = \varphi(a)$,
$S \simeq G/G_s^{=H'}$, $G_s \supset G_a$. L'application $\sigma : A \to A$
correspond à la donnée d'un élément \struck{de}
$\dot{\sigma} \in \mathrm{Norm}(G_a)/G_a$, $\dot{\sigma}$ d'ordre 2 ; le fait
que $\varphi$ \struck{\ill{}} est compatible avec $\sigma$ signifie que
$\sigma \in G_s = H'$, le fait que $\sigma$ opère sans pts fixes signifie que
$\sigma \notin H$ \struck{\ill{}} ; donc on a
\[
\left\{
\begin{array}{l}
\boxed{G \supset H' \supset H} \quad \text{et} \\
\boxed{\dot{\sigma} \in (H' \cap \mathrm{Norm}_G(H))/H \qquad \dot{\sigma}^2 = 1,\ \dot{\sigma} \neq 1}
\end{array}
\right.
\]
\note{entre les deux lignes de l'accolade, une ligne biffée :
\struck{$\sigma \in H' \cap \mathrm{Norm}_G(H) \cap \complement(\ldots)\,H$}, dont
l'argument du complémentaire est noyé sous les ratures}

Le fait que $G$ opère fidèlement sur $\Gamma$ signifie que
\[
\bigcap_{g \in G} g H g^{-1} = \{1\}
\]

Si on a une telle situation ]

\page{33}

\marginal{il existe des arêtes non circulaires}
c) \uncertain{Ou bien} il y a des arêtes non circulaires, \uncertain{ou alors}
les arêtes sont non circulaires, on trouve alors
\[
\underset{a}{A} \xrightarrow[\text{source}]{\ \varphi_1\ } \underset{s}{S}, \qquad
\left\{
\begin{array}{l}
\sigma : A \to A \\
\sigma^2 = \mathrm{id},\ \sigma \text{ ss pt fixe}
\end{array}
\right.
\]
\note{sous $A$ et $S$, « $\ni a$ », « $\ni s$ »}

\[
\boxed{G \supset \underset{G_a}{H'} \supset \underset{G_s}{H} \qquad
\dot{\sigma} \in \mathrm{Norm}_G(H)/H \qquad \dot{\sigma}^2 = 1 \qquad
\dot{\sigma} \neq 1}
\]
\note{sous $H'$ et $H$, des signes d'égalité verticaux renvoient à $G_a$ et
$G_s$. C'est l'inverse de la page~32, où $A \simeq G/G_a$ avec $G_a = H$ et
$S \simeq G/G_s$ avec $G_s = H'$ ; on transcrit tel qu'écrit}

Si on veut que $G$ opère fidèlement, on \struck{\ill{}} doit avoir
\[
\boxed{\bigcap_{g \in G} g^{-1} H g = \{1\}}
\]

Exprimons la condition que $\forall\ \varphi : A \to S \times S$ est
injectif [donc que $\Gamma$ est un « \uncertain{schéma} simplicial de dim.
1 », donné par une \uncertain{ensemble} de parties de cardinal 2 de l'ens.
$S$] i.e. que l'application
\[
\begin{aligned}
G/H &\longrightarrow G/H' \times G/H' \\
gH &\longmapsto gH',\ g\sigma H'
\end{aligned}
\]
donc à prouver que si $g, g' \in G$ tels que $g'^{-1}g \in H'$,
$\sigma^{-1}(g'^{-1}g)\sigma \in H'$, alors \struck{\ill{}} $g'^{-1}g \in H$,
i.e.
\[
\boxed{H' \cap \dot{\sigma}^{-1} H' \dot{\sigma} \subset H}
\]
\note{le $\forall$ devant $\varphi$ est surchargé ; la phrase attend
« … est injective » après l'application, qui n'est pas écrit}

\note{le reste de la page est blanc}


\end{document}
